% Generated draft: requires mathematical and source review.
\documentclass[11pt]{article}
\usepackage[margin=25mm]{geometry}
\usepackage{fontspec}
\setmainfont{Latin Modern Roman}
\usepackage{amsmath,amssymb,mathtools,unicode-math}
\setmathfont{Latin Modern Math}
\usepackage{enumitem,needspace,xurl,hyperref}
\setlength{\parindent}{0pt}
\setlength{\parskip}{5pt}
\setlength{\emergencystretch}{3em}
\newcommand{\N}{\mathbb{N}}
\newcommand{\Z}{\mathbb{Z}}
\newcommand{\Q}{\mathbb{Q}}
\newcommand{\R}{\mathbb{R}}
\newcommand{\C}{\mathbb{C}}
\newcommand{\F}{\mathbb{F}}
\newcommand{\A}{\mathbb{A}}
\newcommand{\PP}{\mathbb{P}}
\newcommand{\E}{\mathbb{E}}
\newcommand{\eps}{\varepsilon}
\newcommand{\dd}{\,\mathrm{d}}
\begin{document}
\section*{UnsolvedMath ID: 30001570}
\textbf{Source problem code:} catalogue entry
\subsection{Definitions and mathematical statement}
The supplied record calls this the global monodromy conjecture for Calabi--Yau varieties. It does not specify the hypotheses on the field \(K\), the convention for a Calabi--Yau variety, or the definition and coefficient ring of the motivic zeta function. The following is therefore a conditional transcription of the supplied claim, pending specification of those data.

For a Calabi--Yau variety \(X/K\) in the intended setting, write \(Z_X(T)\) for its motivic zeta function and \(\mathbb L\) for the Lefschetz class, namely the class of the affine line in the intended motivic coefficient ring \(R\). Let \(K^{\mathrm{sep}}\) be a separable closure of \(K\). Once an admissible prime \(\ell\) and the intended tame monodromy operator \(\sigma\) have been specified, put
\[
H^i_\ell(X)=H^i_{\mathrm{\acute et}}(X\times_K K^{\mathrm{sep}},\Q_\ell),\qquad
\chi_{i,\sigma}(U)=\det\bigl(U\,\mathrm{id}-\sigma\mid H^i_\ell(X)\bigr).
\]
The existence and meaning of this operator require hypotheses absent from the supplied record.

Using the conventional exponent ranges \(a\in\Z\) and \(b\in\Z_{>0}\), define
\[
m(a,b)=\operatorname{ord}_{\Q/\Z}(a/b+\Z)
=\frac{b}{\gcd(|a|,b)}.
\]
Thus \(m(0,b)=1\). Let \(\Phi_m(U)\in\Z[U]\) be the \(m\)-th cyclotomic polynomial.

The supplied conjectural assertion is that, for every such \(X/K\), there exist a finite set \(S\subset\Z\times\Z_{>0}\), positive integers \(e_{a,b}\) for \((a,b)\in S\), and a polynomial \(P(T)\in R[T]\) such that
\[
Z_X(T)=\frac{P(T)}{\displaystyle\prod_{(a,b)\in S}(1-\mathbb L^aT^b)^{e_{a,b}}},
\]
and
\[
\forall(a,b)\in S\quad\exists i\in\Z_{\geq0}\quad
\Phi_{m(a,b)}(U)\mid\chi_{i,\sigma}(U)
\quad\text{in }\Q_\ell[U].
\]
The cohomological degree may depend on \((a,b)\). This asks for the existence of a rational expression with the stated property. The ring \(R\), the interpretation of the equality, the availability of negative powers of \(\mathbb L\), and the exponent ranges must be confirmed from the intended source context; the ranges displayed above are conventional and are not explicitly given in the supplied record. No precise definition of \(Z_X(T)\), or additional hypothesis on \(K\) or \(X\), is supplied here.
\subsection{Short English statement}
The conjecture asks whether every Calabi--Yau variety in the intended setting has a motivic zeta function admitting a finite rational denominator whose exponent ratios predict cyclotomic factors in some tame monodromy characteristic polynomial. Essential field, zeta-function, and monodromy conventions are missing, so this is a conditional statement requiring human mathematical review.
\subsection{Sources}
\begin{itemize}
\item Local UnsolvedMath registry: supplied complete record for permanent ID \(30001570\), source problem code \texttt{catalogue entry}, titled \emph{The Global Monodromy Conjecture for Calabi--Yau Varieties}, in the Oberwolfach Reports Open Problems collection. No individual proposer is supplied; the recorded proposed year is 2010.
\item Supplied source citation: \emph{Komplexe Analysis (2011)}, Oberwolfach Reports, pp.\ 2283--2333, DOI \texttt{10.4171/owr/2010/39}. Original OWR report: \url{https://doi.org/10.4171/owr/2010/39}. The citation title uses 2011, whereas the DOI and recorded proposed year use 2010; their relationship has not been independently checked.
\item Supplied literature reference: Jaspers, global monodromy for K3 surfaces, \url{https://numdam.org/articles/10.1016/j.crma.2016.12.002/}. No fuller bibliographic details are supplied.
\item Literature assessment dated exactly \texttt{2026-08-21}: the supplied curation labels the problem \emph{partially solved} and states that the global monodromy conjecture is proved for tamely ramified abelian varieties and several K3 degenerations, while the general Calabi--Yau statement remains unresolved. This is supplied dated curation derived from the OWR collection, not an independently verified literature assessment; its time-sensitive status and cited claims require independent verification.
\item The registry labels the statement \texttt{corrected\_verified} and reports that its mathematical content was checked against the official Oberwolfach report, with notation, context, and question form normalized without changing the source-level claim. This is a supplied source-verification claim, not verification performed for this draft.
\end{itemize}
\end{document}
